% GENERATED FILE. Edit canonical lessons, then run npm run generate:books.
% canonical-source-hash: fnv1a64:6ed1289fd5905d68
% canonical-lessons: ES-C515-rumor, ES-C515-comunicado, ES-C515-declaracion, ES-C515-declarar, ES-C515-confirmar

\chapter{Según dicen}
\label{ch:es-segun-dicen}
\hlchaptermodality{515}

\begin{chapteropening}
\textbf{By the end of this chapter:} \emph{I can say that there was a rumour, a statement or an official note, and that someone declared or confirmed it.}

Complete the last lesson of chapter 515: I can say that there was a rumour, a statement or an official note, and that someone declared or confirmed it.
\end{chapteropening}

\section[el rumor]{el rumor — the rumour}

\label{lesson:ES-C515-rumor}



\begin{quote}
Before the new one: say the Spanish for to spread, then the Spanish for to announce.
\end{quote}

\subsection*{You'll want to know: el rumor}
\textbf{el rumor} — \textquotedblleft{}the rumour\textquotedblright{}.

\begin{quote}
Corre el rumor de que van a cerrar la fábrica.
\end{quote}

\begin{quote}
\emph{There is a rumour going round that they are going to close the factory.}
\end{quote}

\emph{Corre el rumor} — word goes round. It also means a soft noise: \emph{el rumor del mar}.

\begin{cousinweb}
\emph{rumor} is a learned borrowing from Latin \emph{rūmor} 'common talk, a noise' (EtymDB). The Spanish word kept both senses, the gossip and the murmur.
\end{cousinweb}

\begin{grammarlens}[title={What you've built}]
One more word for telling, explaining and describing at B1.
\end{grammarlens}

\subsection*{Your turn}
\begin{itemize}
\raggedright
  \item \emph{Say it:} \emph{el rumor}
  \item \emph{Say it:} \emph{el rumor}, once more
  \item \emph{Say it:} a sentence of your own with \emph{el rumor}
  \item \emph{Recall:} say the Spanish for to spread, then the Spanish for to announce, then say \emph{el rumor} again
\end{itemize}

\begin{tcolorbox}[breakable,colback=teal!4,colframe=teal!35!black,title={Before you move on}]
What does el rumor mean? (\textquotedblleft{}The rumour\textquotedblright{}.) Say it once more.
\end{tcolorbox}
\section[el comunicado]{el comunicado — the official statement, the press release}

\label{lesson:ES-C515-comunicado}



\begin{quote}
Before the new one: say the Spanish for to announce, then the Spanish for the rumour.
\end{quote}

\subsection*{You'll want to know: el comunicado}
\textbf{el comunicado} — \textquotedblleft{}the official statement, the press release\textquotedblright{}.

\begin{quote}
La empresa envió un comunicado a la prensa.
\end{quote}

\begin{quote}
\emph{The company sent a statement to the press.}
\end{quote}

\emph{Según el comunicado} — according to the statement.

\begin{cousinweb}
\emph{comunicado} is built on the participle of \emph{comunicar} 'to communicate', which comes from Latin \emph{commūnicāre} 'to share' (Lewis \& Short), from \emph{commūnis} 'common'.
\end{cousinweb}

\begin{grammarlens}[title={What you've built}]
One more word for telling, explaining and describing at B1.
\end{grammarlens}

\subsection*{Your turn}
\begin{itemize}
\raggedright
  \item \emph{Say it:} \emph{el comunicado}
  \item \emph{Say it:} \emph{el comunicado}, once more
  \item \emph{Say it:} a sentence of your own with \emph{el comunicado}
  \item \emph{Recall:} say the Spanish for to announce, then the Spanish for the rumour, then say \emph{el comunicado} again
\end{itemize}

\begin{tcolorbox}[breakable,colback=teal!4,colframe=teal!35!black,title={Before you move on}]
What does el comunicado mean? (\textquotedblleft{}The official statement, the press release\textquotedblright{}.) Say it once more.
\end{tcolorbox}
\section[la declaración]{la declaración — the statement; the declaration}

\label{lesson:ES-C515-declaracion}



\begin{quote}
Before the new one: say the Spanish for the rumour, then the Spanish for the official statement.
\end{quote}

\subsection*{You'll want to know: la declaración}
\textbf{la declaración} — \textquotedblleft{}the statement; the declaration\textquotedblright{}.

\begin{quote}
La vecina hizo una declaración a la policía.
\end{quote}

\begin{quote}
\emph{The neighbour made a statement to the police.}
\end{quote}

\emph{Prestar declaración} is to give evidence.

\begin{cousinweb}
\emph{declaración} is a learned borrowing from Latin \emph{dēclārātiō} 'a making clear' (EtymDB), built on \emph{dēclārāre}, from \emph{clārus} 'clear, bright'.
\end{cousinweb}

\begin{grammarlens}[title={What you've built}]
One more word for telling, explaining and describing at B1.
\end{grammarlens}

\subsection*{Your turn}
\begin{itemize}
\raggedright
  \item \emph{Say it:} \emph{la declaración}
  \item \emph{Say it:} \emph{la declaración}, once more
  \item \emph{Say it:} a sentence of your own with \emph{la declaración}
  \item \emph{Recall:} say the Spanish for the rumour, then the Spanish for the official statement, then say \emph{la declaración} again
\end{itemize}

\begin{tcolorbox}[breakable,colback=teal!4,colframe=teal!35!black,title={Before you move on}]
What does la declaración mean? (\textquotedblleft{}The statement; the declaration\textquotedblright{}.) Say it once more.
\end{tcolorbox}
\section[declarar]{declarar — to declare, to state}

\label{lesson:ES-C515-declarar}



\begin{quote}
Before the new one: say the Spanish for the official statement, then the Spanish for the statement.
\end{quote}

\subsection*{You'll want to know: declarar}
\textbf{declarar} — \textquotedblleft{}to declare, to state\textquotedblright{}.

\begin{quote}
El presidente declaró que no sabía nada.
\end{quote}

\begin{quote}
\emph{The president stated that he knew nothing.}
\end{quote}

In court it means to testify, and at customs to declare goods: \emph{¿Algo que declarar?}

\begin{cousinweb}
\emph{declarar} comes from Latin \emph{dēclārāre} 'to make clear, to show' (EtymDB), a learned word built on \emph{clārus} 'clear'. English \emph{declare} is the same Latin verb.
\end{cousinweb}

\begin{grammarlens}[title={What you've built}]
One more word for telling, explaining and describing at B1.
\end{grammarlens}

\subsection*{Your turn}
\begin{itemize}
\raggedright
  \item \emph{Say it:} \emph{declarar}
  \item \emph{Say it:} \emph{declarar}, once more
  \item \emph{Say it:} a sentence of your own with \emph{declarar}
  \item \emph{Recall:} say the Spanish for the official statement, then the Spanish for the statement, then say \emph{declarar} again
\end{itemize}

\begin{tcolorbox}[breakable,colback=teal!4,colframe=teal!35!black,title={Before you move on}]
What does declarar mean? (\textquotedblleft{}To declare, to state\textquotedblright{}.) Say it once more.
\end{tcolorbox}
\section[confirmar]{confirmar — to confirm}

\label{lesson:ES-C515-confirmar}



\begin{quote}
Before the new one: say the Spanish for the statement, then the Spanish for to declare.
\end{quote}

\subsection*{You'll want to know: confirmar}
\textbf{confirmar} — \textquotedblleft{}to confirm\textquotedblright{}.

\begin{quote}
La policía confirmó la noticia.
\end{quote}

\begin{quote}
\emph{The police confirmed the news.}
\end{quote}

\emph{Confirmar una reserva} — to confirm a booking.

\begin{cousinweb}
\emph{confirmar} is a learned borrowing from Latin \emph{cōnfirmāre} 'to make firm, to strengthen' (EtymDB), built on \emph{con-} and \emph{firmāre}, from \emph{firmus} 'firm'.
\end{cousinweb}

\begin{grammarlens}[title={What you've built}]
One more word for telling, explaining and describing at B1.
\end{grammarlens}

\subsection*{Your turn}
\begin{itemize}
\raggedright
  \item \emph{Say it:} \emph{confirmar}
  \item \emph{Say it:} \emph{confirmar}, once more
  \item \emph{Say it:} a sentence of your own with \emph{confirmar}
  \item \emph{Recall:} say the Spanish for the statement, then the Spanish for to declare, then say \emph{confirmar} again
\end{itemize}

\begin{tcolorbox}[breakable,colback=teal!4,colframe=teal!35!black,title={Before you move on}]
What does confirmar mean? (\textquotedblleft{}To confirm\textquotedblright{}.) Say it once more.
\end{tcolorbox}
